\documentclass[11pt, a4paper]{article}

% Gemeinsame Style-Definitionen (TEXINPUTS muss templates/ enthalten — siehe scripts/build_tex.py)
% bfw-style.tex — gemeinsame Style-Definitionen für alle Handouts/Aufgabenhefte
% Voraussetzung: Kompilation mit xelatex (wegen fontspec)

\usepackage[ngerman]{babel}
\usepackage{geometry}
\geometry{a4paper, top=2.2cm, bottom=2.2cm, left=2.3cm, right=2.3cm}

\usepackage{fontspec}
% Fallback-Kette: Source Sans Pro -> Calibri -> Segoe UI -> Latin Modern Sans
% (Latin Modern ist immer mit MiKTeX vorhanden)
\IfFontExistsTF{Source Sans Pro}{%
  \setmainfont{Source Sans Pro}[Ligatures=TeX]%
}{\IfFontExistsTF{Calibri}{%
  \setmainfont{Calibri}[Ligatures=TeX]%
}{\IfFontExistsTF{Segoe UI}{%
  \setmainfont{Segoe UI}[Ligatures=TeX]%
}{%
  \setmainfont{Latin Modern Sans}[Ligatures=TeX]%
}}}
\IfFontExistsTF{Source Code Pro}{%
  \setmonofont{Source Code Pro}[Scale=0.92]%
}{\IfFontExistsTF{Consolas}{%
  \setmonofont{Consolas}[Scale=0.92]%
}{%
  \setmonofont{Latin Modern Mono}[Scale=0.92]%
}}

\usepackage{xcolor}
% Farbpalette: hell, freundlich, modern
\definecolor{bfwPrimary}{HTML}{0EA5A4}    % Petrol/Türkis - Primär
\definecolor{bfwPrimaryDark}{HTML}{0F766E} % dunkler Petrol für Headings
\definecolor{bfwAccent}{HTML}{F59E0B}     % Amber - Akzent
\definecolor{bfwSuccess}{HTML}{10B981}    % Grün - Erfolg / Lösung
\definecolor{bfwWarn}{HTML}{F97316}       % Orange - Achtung
\definecolor{bfwInk}{HTML}{0F172A}        % fast Schwarz
\definecolor{bfwInkSoft}{HTML}{475569}    % gedämpftes Grau
\definecolor{bfwBgSoft}{HTML}{F8FAFC}     % off-white
\definecolor{bfwBgTeal}{HTML}{ECFEFF}     % helles Türkis
\definecolor{bfwBgAmber}{HTML}{FEF3C7}    % helles Gelb
\definecolor{bfwBgRose}{HTML}{FFE4E6}     % helles Rosé für Eselsbrücken

\usepackage[colorlinks=true, linkcolor=bfwPrimaryDark, urlcolor=bfwPrimaryDark]{hyperref}
\usepackage{enumitem}
\setlist[itemize]{leftmargin=*, itemsep=2pt, topsep=2pt}
\setlist[enumerate]{leftmargin=*, itemsep=2pt, topsep=2pt}

\usepackage[most]{tcolorbox}
\usepackage{booktabs}
\usepackage{tikz}
\usetikzlibrary{decorations.pathmorphing, positioning, shapes.geometric, arrows.meta}

% Merkkästchen — wichtige Definition / Kernpunkt
\newtcolorbox{merksatz}[1][]{%
  enhanced, breakable,
  colback=bfwBgTeal,
  colframe=bfwPrimary,
  boxrule=0.6pt,
  arc=4pt,
  left=10pt, right=10pt, top=8pt, bottom=8pt,
  fonttitle=\bfseries\color{bfwPrimaryDark},
  title={\faLightbulb\ Merksatz},
  #1
}

% Eselsbrücke — Mnemoniks
\newtcolorbox{eselsbruecke}[1][]{%
  enhanced, breakable,
  colback=bfwBgRose,
  colframe=bfwAccent,
  boxrule=0.6pt,
  arc=4pt,
  left=10pt, right=10pt, top=8pt, bottom=8pt,
  fonttitle=\bfseries\color{bfwWarn},
  title={Eselsbrücke},
  #1
}

% Beispiel-Box
\newtcolorbox{beispiel}[1][]{%
  enhanced, breakable,
  colback=bfwBgSoft,
  colframe=bfwInkSoft,
  boxrule=0.4pt,
  arc=3pt,
  left=10pt, right=10pt, top=8pt, bottom=8pt,
  fonttitle=\bfseries\color{bfwInk},
  title={Beispiel},
  #1
}

% Lösungs-Box (für Aufgabenhefte)
\newtcolorbox{loesung}[1][]{%
  enhanced, breakable,
  colback=bfwBgAmber!40,
  colframe=bfwSuccess,
  boxrule=0.6pt,
  arc=3pt,
  left=10pt, right=10pt, top=8pt, bottom=8pt,
  fonttitle=\bfseries\color{bfwSuccess!70!black},
  title={Lösung},
  #1
}

% Glossar-Eintrag
\newcommand{\glossareintrag}[2]{%
  \par\noindent\textbf{\color{bfwPrimaryDark}#1} \enspace #2\par\smallskip
}

% Section-Stil — etwas farbiger als Standard
\usepackage{titlesec}
\titleformat{\section}
  {\Large\bfseries\color{bfwPrimaryDark}}
  {\thesection}{0.6em}{}
\titleformat{\subsection}
  {\large\bfseries\color{bfwPrimary}}
  {\thesubsection}{0.5em}{}
\titleformat{\subsubsection}
  {\normalsize\bfseries\color{bfwInk}}
  {\thesubsubsection}{0.4em}{}

% TikZ-Stil "skizzenhaft" — etwas weicher als Rough.js, aber stabil
\tikzset{
  roughbox/.style = {
    draw=bfwPrimaryDark, thick, fill=bfwBgTeal,
    rounded corners=4pt, inner sep=6pt
  },
  roughaccent/.style = {
    draw=bfwAccent, thick, fill=bfwBgAmber,
    rounded corners=4pt, inner sep=6pt
  },
  roughline/.style = {
    draw=bfwInkSoft, thick, -{Stealth[length=2.5mm]}
  }
}

% Bullet-Symbole (bewusst kein FontAwesome — vermeidet Font-Installations-Probleme)
\newcommand{\faLightbulb}{$\bullet$}
\newcommand{\faBookmark}{$\bullet$}

% Header/Footer
\usepackage{fancyhdr}
\pagestyle{fancy}
\fancyhf{}
\renewcommand{\headrulewidth}{0pt}
\fancyfoot[L]{\small\color{bfwInkSoft}\thepage}
\fancyfoot[R]{\small\color{bfwInkSoft} BFW M-064 Beschaffung im Büromanagement}


\title{\vspace*{-1.5cm}\Huge\bfseries\color{bfwPrimaryDark}Tag~5: Zeitmanagement \& Aufgabenplanung\\[0.4em]
  \large\color{bfwInkSoft}\textnormal{Die eigene Zeit bewusst steuern}}
\author{\textsf{Büroprozesse gestalten und Arbeitsvorgänge organisieren \textperiodcentered{} M-067}\\
\textsf{\small\copyright\ Can Siebert\,\textperiodcentered\,2026}}
\date{03.07.2026}

\begin{document}
\maketitle
\thispagestyle{empty}

\vspace{1em}
\begin{merksatz}[title={\bfseries Worum geht es heute?}]
24~Stunden sind 24~Stunden --- jeder Mensch hat gleich viel Zeit am Tag. Den Unterschied
macht, \emph{wie} wir sie nutzen. Beim Zeitmanagement geht es darum, die persönlich
vorhandene Zeit möglichst effektiv einzusetzen: Stress und Überforderung vermeiden,
Wichtiges von Unwichtigem trennen, Prioritäten setzen und Beruf und Privatleben in
Balance halten. In diesem Handout lernen Sie, wie Sie Zeitdiebe erkennen, Ihre
Leistungskurve nutzen und mit bewährten Methoden --- ALPEN, Eisenhower, Pareto, ABC ---
Ihren Arbeitstag planen. Das ist die Grundlage für Tag~6 (Terminplanung) und für die
IHK-Abschlussprüfung.
\end{merksatz}

\tableofcontents
\vspace{1em}
\hrule
\vspace{1em}

\section{Was Zeitmanagement leisten soll}

24~Stunden hat jeder Tag --- aber jeder Mensch nutzt sie unterschiedlich. Beim
\emph{Zeitmanagement} geht es darum, die persönlich vorhandene Zeit möglichst effektiv
zu nutzen. Dafür gibt es verschiedene Hilfsmittel und Methoden. Für die Planung muss man
zwar etwas Zeit investieren, doch bei guter Planung gewinnt man im Anschluss viel Zeit
für berufliche \emph{und} private Dinge zurück.

Mit dem Zeitmanagement werden mehrere \emph{Ziele} verfolgt: die Vermeidung von Stress
und Überforderung, die Steigerung der persönlichen Zufriedenheit und genügend Zeit für
die wesentlichen Aufgaben in Beruf und Privatleben. Privat- und Berufsleben sollen in
einem ausgewogenen Verhältnis stehen (\emph{Work-Life-Balance}), und die Freizeit soll
ohne Reue genossen werden können.

\begin{merksatz}
Für ein optimales Zeitmanagement ist es notwendig, \textbf{Zeitfresser zu erkennen},
\textbf{wichtige von unwichtigen Aufgaben zu unterscheiden} und \textbf{Prioritäten zu
setzen}.
\end{merksatz}

\section{Störungen --- die Zeitdiebe}

Die Beseitigung von Störungen (Zeitfressern, Zeitdieben, Zeitfallen) im Arbeitsalltag
gehört zu den Hauptaufgaben des Zeitmanagements, denn durch sie geht wertvolle Zeit
verloren. Man unterteilt sie in \emph{äußere} und \emph{innere} Störungen. Innere
Störungen werden durch die Person selbst verursacht, äußere Störungen durch andere
Personen oder durch die äußeren Gegebenheiten.

\begin{center}
\begin{tabular}{p{6.4cm} p{6.4cm}}
\toprule
\textbf{Äußere Störungen} & \textbf{Innere Störungen}\\
\midrule
unordentlicher Arbeitsplatz & Lustlosigkeit\\
Unterbrechung durch Kollegen (Nachfragen) & Unfähigkeit, „nein” zu sagen\\
klingelnde Telefone, ständige E-Mail-Kontrolle & fehlende Selbstdisziplin\\
Lärm (Straßenverkehr, Geräte, Publikum) & Angstgefühle, beunruhigende Gedanken\\
Wartezeiten am Gemeinschaftsdrucker & Redseligkeit\\
technische Störungen (Computerausfall) & Müdigkeit, „vor sich hin träumen”\\
\bottomrule
\end{tabular}
\end{center}

\begin{eselsbruecke}
\textbf{Außen oder innen?} Frage dich: \emph{„Kommt die Störung von außen (anderen
Personen, Technik, Umgebung) --- oder von mir selbst (Stimmung, Disziplin)?”} Innere
Störungen kann man durch Selbstmanagement angehen, äußere durch Organisation und Regeln.
\end{eselsbruecke}

\section{Der Sägeblatteffekt}

Gerade haben Sie sich in eine schwierige Aufgabe eingearbeitet und sind hochkonzentriert
--- da klingelt das Telefon und reißt Sie aus Ihren Gedanken. Nach dem Gespräch brauchen
Sie erst wieder eine gewisse Zeit, um in die Aufgabe zurückzufinden und Ihre
Konzentration auf das alte Maß zu bringen. Kommen solche Störungen häufig vor, entsteht
der sogenannte \emph{Sägeblatteffekt}: Mit jeder Unterbrechung fällt die
Leistungsfähigkeit ab und muss anschließend wieder aufgebaut werden. Stellt man Zeit und
Leistungsfähigkeit grafisch dar, ergibt sich ein Bild, das einem Sägeblatt ähnelt.

\begin{merksatz}
Addiert man alle Leistungsverluste durch Störungen auf, können bis zu \textbf{28~\%} der
Zeit verloren gehen.
\end{merksatz}

Ein Gegenmittel ist die \emph{stille Stunde}: ein fest geplanter Termin mit sich selbst
für ungestörtes Arbeiten an wichtigen Aufgaben. Tragen Sie diese Zeit in Ihren
Terminplan ein, stellen Sie das Telefon auf Kollegen um oder schalten Sie den
Anrufbeantworter ein. Legen Sie die stille Stunde möglichst in einen Teil des Tages mit
ohnehin wenigen Störungen --- und nutzen Sie sie wirklich für wichtige Aufgaben, nicht
etwa zum Kontrollieren des E-Mail-Postfachs.

\section{Die Leistungskurve}

Bei der Verteilung der Aufgaben über den Arbeitstag sollten Sie Ihre persönliche
\emph{Leistungskurve} beachten. Die meisten Menschen haben eine durchschnittliche
Leistungskurve: Das Leistungshoch liegt zwischen \textbf{08:00 und 12:00 Uhr}. Nach der
Mittagspause folgt das größte Tief. Deshalb sollten alle Arbeiten, für die volle
Leistungsfähigkeit nötig ist, am Vormittag erledigt werden; am Nachmittag eignen sich
Routineaufgaben wie die Ablage von Unterlagen.

Wer keine durchschnittliche Leistungskurve hat, ist entweder ein \emph{Morgenmensch}
(höchste Leistung am frühen Vormittag) oder ein \emph{Abendmensch} (höchste Leistung am
Nachmittag). Beide Typen sollten die wichtigsten Aufgaben in ihre persönliche Hochphase
legen.

\subsection{Pausen und das Arbeitszeitgesetz}

Wichtig sind regelmäßige \emph{Pausen}. Sie dienen dazu, Kraft für die nächsten Aufgaben
zu tanken. Studien zeigen: Pausen machen produktiver, senken den Stresspegel, steigern
die Leistungsfähigkeit und verringern das Burnout-Risiko; laut DGUV sinkt zudem das
Risiko für Arbeitsunfälle. Nicht als Pause gilt, wenn man sein Brötchen isst und dabei
weiter E-Mails liest --- gedanklich ist man dann noch bei der Arbeit. Ein paar Schritte
gehen oder die Augen kurz schließen kann dagegen sehr erholsam sein.

\begin{beispiel}
\textbf{§~4 ArbZG (Auszug):} Die Arbeit ist durch im Voraus feststehende Ruhepausen von
mindestens \textbf{30~Minuten} bei einer Arbeitszeit von mehr als sechs bis zu neun
Stunden und \textbf{45~Minuten} bei mehr als neun Stunden zu unterbrechen. Die Pausen
können in Zeitabschnitte von jeweils mindestens 15~Minuten aufgeteilt werden. Länger als
sechs Stunden hintereinander dürfen Arbeitnehmer nicht ohne Ruhepause beschäftigt werden.
\end{beispiel}

\section{Methoden des Zeitmanagements}

Das Zeitmanagement bietet verschiedene Methoden an. Im Folgenden lernen Sie die fünf
wichtigsten kennen: ABC-Analyse, Eisenhower-Prinzip, ALPEN-Methode, Pareto-Prinzip und
Salamitaktik.

\subsection{Die ABC-Analyse}

Bei der ABC-Analyse teilt man die Aufgaben nach ihrer \emph{Priorität} (Wichtigkeit) in
drei Kategorien ein. Dazu müssen die Aufgaben analysiert und gewichtet werden, denn nicht
alle sind gleich bedeutend.

\begin{center}
\begin{tabular}{p{2.4cm} p{3.2cm} p{6.4cm}}
\toprule
\textbf{Kategorie} & \textbf{Priorität} & \textbf{Was zu beachten ist}\\
\midrule
A-Aufgaben & hoch (Muss) & Wichtigste Aufgaben; morgens beginnen, wenn man ausgeruht ist; selbst erledigen, nicht delegieren.\\
B-Aufgaben & mittel (Soll) & Nach den A-Aufgaben erledigen; meist nachmittags.\\
C-Aufgaben & niedrig (Kann) & Routine; benötigen oft die meiste Zeit, tragen aber am wenigsten zum Erfolg bei.\\
\bottomrule
\end{tabular}
\end{center}

\medskip
Typisch ist ein Missverhältnis: C-Aufgaben verschlingen rund 65~\% der Zeit, bringen aber
nur 15~\% des Erfolgs; A-Aufgaben kosten nur 15~\% der Zeit, bringen aber 65~\% des
Erfolgs. Wer im Leistungstief steckt oder es hektisch wird, neigt zu C-Aufgaben --- und
verliert die wirklich wichtigen aus dem Blick.

\subsection{Das Eisenhower-Prinzip}

Benannt nach dem US-General und späteren US-Präsidenten Dwight~D.\ Eisenhower. Aufgaben
werden nach zwei Kriterien sortiert: \emph{Dringlichkeit} und \emph{Wichtigkeit} ---
wobei \textbf{Wichtigkeit vor Dringlichkeit} geht. Dringende Aufgaben haben einen
konkreten Termin; wichtige Aufgaben bringen ein Unternehmen oder eine Person einem Ziel
näher (Arbeitserfolg) und sind oft langfristig (strategisch) orientiert.

\begin{center}
\begin{tabular}{p{2.6cm} p{4.7cm} p{4.7cm}}
\toprule
 & \textbf{dringend} & \textbf{nicht dringend}\\
\midrule
\textbf{wichtig} & \textbf{Feld 1:} sofort und selbst erledigen (Notfälle, Abgabetermine) & \textbf{Feld 2:} Termin vereinbaren, in die Terminplanung aufnehmen\\
\textbf{nicht wichtig} & \textbf{Feld 3:} delegieren (z.\,B.\ Software-Neuinstallation) & \textbf{Feld 4:} Papierkorb\\
\bottomrule
\end{tabular}
\end{center}

\begin{eselsbruecke}
\textbf{„Selbst --- Termin --- Delegieren --- Papierkorb”} für die Felder 1--4. Erst
Wichtigkeit prüfen, dann Dringlichkeit --- denn wichtige Dinge dürfen nicht im
Tagesgeschäft untergehen.
\end{eselsbruecke}

\subsection{Die ALPEN-Methode}

Der Wirtschaftswissenschaftler Lothar~J.\ Seiwert entwickelte die \emph{ALPEN-Methode}.
Dabei nehmen Sie sich jeden Tag ein paar Minuten Zeit, um zu Beginn des Arbeitstages
einen schriftlichen Tagesplan zu erstellen. Die fünf Schritte:

\begin{enumerate}
  \item \textbf{A} --- \textbf{A}ufgaben und Termine zusammenstellen und aufschreiben.
  \item \textbf{L} --- \textbf{L}änge der Tätigkeiten realistisch abschätzen.
  \item \textbf{P} --- \textbf{P}ufferzeiten einplanen.
  \item \textbf{E} --- \textbf{E}ntscheidungen über Prioritäten treffen, eventuell delegieren.
  \item \textbf{N} --- \textbf{N}achkontrolle: alles, was nicht erledigt werden konnte, neu einplanen.
\end{enumerate}

\begin{merksatz}
Verplanen Sie nur ca.\ \textbf{60~\%} Ihrer Zeit. Die restlichen 40~\% bleiben als Reserve
für Unvorhergesehenes --- sonst gerät der Plan schon bei kleinen Abweichungen durcheinander.
\end{merksatz}

\subsection{Das Pareto-Prinzip (80/20-Regel)}

Anfang des 20.\ Jahrhunderts entdeckte der italienische Wirtschaftswissenschaftler
Vilfredo Pareto, dass 20~\% der Bevölkerung 80~\% des Staatsvermögens besaßen. Daraus
entstand die \emph{80/20-Regel}. Auf das Zeitmanagement bezogen heißt das: \textbf{20~\%
der eingesetzten Zeit bewirken 80~\% des Erfolgs}; die restlichen 80~\% der Zeit ergeben
nur noch 20~\% des Erfolgs.

\begin{beispiel}
20~\% der Arbeit am Schreibtisch ermöglichen häufig 80~\% des Arbeitserfolgs.\\
20~\% der Kunden eines Unternehmens sorgen für 80~\% des Gewinns.
\end{beispiel}

Daraus folgt: Es macht wenig Sinn, alle Aufgaben zu 100~\% und mit gleicher Aufmerksamkeit
erledigen zu wollen. Prioritäten setzen und Wichtiges zuerst!

\subsection{Die Salamitaktik}

Große Aufgaben bereiten Schwierigkeiten, weil sie schwer zu planen sind und ein „Berg
Arbeit” demotivierend wirkt. Bei der \emph{Salamitaktik} zerlegt man eine komplexe,
große Aufgabe in kleinere, überschaubare Schritte (sinnbildlich in Salamischeiben) und
erledigt sie Schritt für Schritt. So entstehen immer wieder kleine Erfolgserlebnisse,
und eine zunächst unmöglich erscheinende Aufgabe wird machbar.

\begin{merksatz}
Alle Methoden haben Stärken und Schwächen. Wählen Sie die Methode, mit der Sie persönlich
am besten zurechtkommen --- nur eine richtig angewandte Methode spart wirklich Zeit.
\end{merksatz}

\section{Weitere Hilfsmittel: Listen und Ziele}

\subsection{Checkliste und To-do-Liste}

Eine \emph{Checkliste} koordiniert Aktivitäten innerhalb eines Prozesses und überprüft
sie auf Vollständigkeit und Reihenfolge. Sie unterstützt systematisches, fehlerfreies
Arbeiten, erleichtert die Einarbeitung neuer Kollegen und ist
\emph{vergangenheitsbezogen} (Wurde alles erledigt?).

Eine \emph{To-do-Liste} (Aufgabenliste) hält dagegen zukünftige Aufgaben fest und
organisiert sie. Dahinter steckt die Frage: \emph{wer macht was bis wann?} Sie ist durch
ihre Fragestellung „Was ist zu tun?” auf die Zukunft gerichtet ---
\emph{zukunftsbezogen}.

\begin{beispiel}
\textbf{Modern umgesetzt:} Digitale Aufgabenlisten wie \emph{Microsoft To~Do},
\emph{Todoist} oder \emph{Trello} verbinden To-do-Liste, Prioritäten und Delegation:
Erinnerungen, Fälligkeitsdaten, Zuweisung an Kollegen, jederzeit änder- und teilbar. Die
\emph{Pomodoro-Technik} (z.\,B.\ 25~Min konzentriert arbeiten, dann 5~Min Pause) schützt
vor dem Sägeblatteffekt --- eine moderne Form der stillen Stunde.
\end{beispiel}

\subsection{Ziele formulieren --- die SMART-Regel}

Wirksame Ziele zu formulieren ist im Arbeitsalltag wichtig. Die \emph{SMART-Regel} nennt
fünf Kriterien:

\begin{itemize}
  \item \textbf{S}pezifisch --- so präzise wie möglich, nicht vage.
  \item \textbf{M}essbar --- die Zielerreichung muss sich messen lassen.
  \item \textbf{A}kzeptiert --- das Ziel wird von den Betroffenen angenommen (positiv formulieren).
  \item \textbf{R}ealistisch --- das Ziel muss erreichbar sein.
  \item \textbf{T}erminiert --- mit genauer Vorgabe, wann es erreicht sein muss.
\end{itemize}

\begin{beispiel}
„Verringerung des Druckerpapierverbrauchs in der Abteilung Beschaffung der Duisdorfer
Büro-Konzept~KG um 10~\% bis zum 31.12.” --- spezifisch, messbar, terminiert.
\end{beispiel}

\section{Selbstmanagement}

Beim \emph{Selbstmanagement} übernimmt jeder die Verantwortung für die eigene
Entwicklung --- persönlich wie beruflich. Das Verhalten beim Lernen oder bei der Arbeit
wird so gesteuert und strukturiert, dass die gesetzten Ziele erreicht werden können.
Selbstmanagement wird durch den Willen der Person gelenkt und durch Vorlieben,
Abneigungen und Erfahrungen beeinflusst. Wichtig ist, nicht nach gut/schlecht zu
urteilen, sondern zu hinterfragen, welche Strategie das angestrebte Ziel (z.\,B.\ die
IHK-Abschlussprüfung) unterstützt und welche eher hinderlich ist.

\begin{merksatz}
Teilkompetenzen des Selbstmanagements: \textbf{sich selbst motivieren, sich Ziele setzen,
planen, sich selbst organisieren, sich Feedback geben, den eigenen Erfolg überprüfen, die
eigene Lernfähigkeit nutzen}.
\end{merksatz}

\section{Zusammenfassung}

\begin{enumerate}
  \item Zeitmanagement nutzt die vorhandene Zeit effektiv: \textbf{Stress vermeiden, Prioritäten setzen, Work-Life-Balance}.
  \item \textbf{Störungen/Zeitdiebe} sind äußere (von anderen/Umgebung) oder innere (von einem selbst).
  \item Der \textbf{Sägeblatteffekt} kostet bis zu 28~\% der Zeit; Gegenmittel ist die \textbf{stille Stunde}.
  \item Die \textbf{Leistungskurve} legt das Hoch auf 08--12~Uhr; wichtige Aufgaben gehören in die persönliche Hochphase. Pausen nach \textbf{§~4 ArbZG}.
  \item Fünf Methoden: \textbf{ABC-Analyse, Eisenhower-Prinzip, ALPEN-Methode, Pareto-Prinzip, Salamitaktik}.
  \item Hilfsmittel: \textbf{Checkliste} (Vergangenheit), \textbf{To-do-Liste} (Zukunft), \textbf{SMART-Ziele} --- getragen vom \textbf{Selbstmanagement}.
\end{enumerate}

\newpage

\section*{Glossar}
\addcontentsline{toc}{section}{Glossar}

\glossareintrag{Zeitmanagement}{Effektive Nutzung der persönlich vorhandenen Zeit mit Hilfsmitteln und Methoden; Ziele sind Stressvermeidung, Zufriedenheit und Work-Life-Balance.}

\glossareintrag{Zeitdieb / Störung}{Zeitfresser im Arbeitsalltag; unterteilt in äußere (von anderen/Umgebung) und innere (von der Person selbst) Störungen.}

\glossareintrag{Sägeblatteffekt}{Leistungsverlust durch häufige Unterbrechungen; der Kurvenverlauf ähnelt einem Sägeblatt. Bis zu 28~\% Zeitverlust.}

\glossareintrag{Stille Stunde}{Fest geplanter Termin mit sich selbst für ungestörtes Arbeiten an wichtigen Aufgaben.}

\glossareintrag{Leistungskurve}{Verlauf der Leistungsfähigkeit über den Tag; Hoch meist 08--12~Uhr, Tief nach der Mittagspause.}

\glossareintrag{Morgen-/Abendmensch}{Personen mit höchster Leistungsfähigkeit am frühen Vormittag bzw.\ am Nachmittag.}

\glossareintrag{§~4 ArbZG}{Regelt Ruhepausen: mind.\ 30~Min bei >6--9~h, 45~Min bei >9~h; aufteilbar in 15-Min-Abschnitte.}

\glossareintrag{ABC-Analyse}{Einteilung der Aufgaben nach Priorität: A = Muss (hoch), B = Soll (mittel), C = Kann (niedrig).}

\glossareintrag{Eisenhower-Prinzip}{Sortierung nach Dringlichkeit und Wichtigkeit; Wichtigkeit vor Dringlichkeit. Felder: selbst erledigen, Termin, delegieren, Papierkorb.}

\glossareintrag{ALPEN-Methode}{Tagesplan in fünf Schritten: Aufgaben aufschreiben, Länge schätzen, Pufferzeiten, Entscheidungen/delegieren, Nachkontrolle. Nur ca.\ 60~\% verplanen.}

\glossareintrag{Pareto-Prinzip}{80/20-Regel: 20~\% der Zeit bewirken 80~\% des Erfolgs.}

\glossareintrag{Salamitaktik}{Zerlegen großer Aufgaben in kleine, überschaubare Schritte (Salamischeiben).}

\glossareintrag{Checkliste}{Liste zur Prüfung von Vollständigkeit und Reihenfolge; vergangenheitsbezogen.}

\glossareintrag{To-do-Liste}{Aufgabenliste für zukünftige Aufgaben (wer macht was bis wann); zukunftsbezogen.}

\glossareintrag{SMART-Regel}{Kriterien für Ziele: Spezifisch, Messbar, Akzeptiert, Realistisch, Terminiert.}

\glossareintrag{Selbstmanagement}{Eigenverantwortliche Steuerung des eigenen Verhaltens zur Erreichung gesetzter Ziele.}

\glossareintrag{Pomodoro-Technik}{Moderne Methode: Wechsel aus festen Arbeitsintervallen (z.\,B.\ 25~Min) und kurzen Pausen.}

\end{document}
